% Generated draft: requires mathematical and source review.
\documentclass[11pt]{article}
\usepackage[margin=25mm]{geometry}
\usepackage{fontspec}
\setmainfont{Latin Modern Roman}
\usepackage{amsmath,amssymb,mathtools,unicode-math}
\setmathfont{Latin Modern Math}
\usepackage{enumitem,needspace,xurl,hyperref}
\setlength{\parindent}{0pt}
\setlength{\parskip}{5pt}
\setlength{\emergencystretch}{3em}
\newcommand{\N}{\mathbb{N}}
\newcommand{\Z}{\mathbb{Z}}
\newcommand{\Q}{\mathbb{Q}}
\newcommand{\R}{\mathbb{R}}
\newcommand{\C}{\mathbb{C}}
\newcommand{\F}{\mathbb{F}}
\newcommand{\A}{\mathbb{A}}
\newcommand{\PP}{\mathbb{P}}
\newcommand{\E}{\mathbb{E}}
\newcommand{\eps}{\varepsilon}
\newcommand{\dd}{\,\mathrm{d}}
\begin{document}
\section*{UnsolvedMath ID: 700357}
\textbf{Source problem code:} catalogue entry
\subsection{Definitions and mathematical statement}
Let \(p\) be a prime and let \(n\geq 1\) be an integer satisfying \(p>2n+1\). Here positivity of \(n\) is the usual matrix-size convention for \(\mathrm{GL}_n(\C)\); the source states only that \(n\) is an integer. Write \(\mathrm{GL}_n(\C)\) for the group of invertible \(n\times n\) complex matrices. A \(p\)-element is an element of finite order equal to a power of \(p\), with the identity allowed. Let \(x,y\in\mathrm{GL}_n(\C)\) be \(p\)-elements, so that
\[
\operatorname{ord}(x)=p^a,\qquad \operatorname{ord}(y)=p^b
\quad\text{for some }a,b\in\Z_{\geq 0},
\]
and put \(H=\langle x,y\rangle\), the subgroup generated by \(x\) and \(y\).

The source attributes to W. Feit and J. G. Thompson the following finite-case result: if \(H\) is finite, then \(H\) is an abelian \(p\)-group, meaning that it is abelian and has order \(p^c\) for some \(c\in\Z_{\geq 0}\).

\textbf{Problem (J. D. Dixon, 1978).} Under the same hypotheses on \(p,n,x,y\), assume instead that \(H\) is infinite. What structural conclusions can be drawn about \(H\)?

The supplied formulation is an open-ended request for structural information. It does not specify a particular property to prove, a classification to obtain, or a proposed analogue of the finite-case conclusion. Thus no more specific conjectural assertion is determined by the record.
\subsection{Short English statement}
J. D. Dixon asks what can be said about an infinite subgroup of \(\mathrm{GL}_n(\C)\) generated by two elements of finite \(p\)-power order when \(p>2n+1\). The source contrasts this with an attributed theorem that every such finite subgroup is an abelian \(p\)-group. No particular conclusion for the infinite case is specified.
\subsection{Sources}
\begin{itemize}
\item Local UnsolvedMath registry, permanent ID \(700357\), exact source problem code \texttt{catalogue entry}, title \emph{Kourovka Notebook Problem 6.10}. The supplied registry labels the problem \texttt{open}; this is supplied status information, not an independently verified assessment.
\item \emph{The Kourovka Notebook (Unsolved Problems in Group Theory)}, original problem 6.10, 6th issue (1978), proposed by J. D. Dixon. The supplied background identifies the consulted edition as the 21st issue edition (2026), arXiv:1401.0300v48. The supplied collection description names E. I. Khukhro and V. D. Mazurov as editors. These bibliographic details have not been independently verified here.
\item W. Feit and J. G. Thompson, \emph{Pacif. J. Math.}, volume 11, no. 4 (1961), pp. 1257--1262, cited in the supplied statement for the finite-case result. No article title is supplied.
\end{itemize}
\end{document}
